\section{An exponential--rational hypergeometric integral}
\label{app:hypergeometric-integral}

The geometric-series derivation of \(\Psi\) in
\cref{sec:transfer-birth-death} is the shortest route to the generating function.
Direct integration of the rational characteristic leads to the following more
general identity.

\begin{proposition}
\label{prop:hypergeometric-integral}
Let \(h,k,p,q\) be real constants with \(hk\neq0\), \(p\neq0\) and
\(h/k>0\).  On any interval on which the denominator is non-zero,
\begin{equation}
  \int
  \frac{\e^{h\tau}}{p+q\e^{k\tau}}\,\dd\tau
  =
  \frac{\e^{h\tau}}{ph}\,
  {}_2F_1\!\left(
    1,\frac hk;\frac hk+1;
    -\frac qp\e^{k\tau}
  \right)+C,
  \label{eq:hypergeometric-integral}
\end{equation}
where the hypergeometric function is continued from its principal power series
along the interval.  In particular, the Euler-integral proof below applies
directly while
\(-q\e^{k\tau}/p\notin[1,\infty)\).
\end{proposition}

\begin{proof}
Use \(v=\e^{k\tau}\), so that
\(\dd\tau=\dd v/(kv)\).  Then
\begin{equation}
  \int\frac{\e^{h\tau}}{p+q\e^{k\tau}}\,\dd\tau
  =\frac1{pk}
   \int\frac{v^{h/k-1}}{1+(q/p)v}\,\dd v.
  \label{eq:hyp-proof-substitution}
\end{equation}
Write \(a=h/k>0\).  An antiderivative may be fixed, up to a constant, by
integrating from \(0\) to \(v\):
\begin{equation}
  \frac1{pk}\int_0^v
  \frac{w^{a-1}}{1+(q/p)w}\,\dd w.
  \label{eq:hyp-proof-definite}
\end{equation}
The lower endpoint is integrable precisely because \(a>0\).  Rescale
\(w=vu\) to obtain
\begin{equation}
  \frac{v^a}{pk}
  \int_0^1
  u^{a-1}
  \{1-zu\}^{-1}\,\dd u,
  \qquad
  z=-\frac qp v.
  \label{eq:hyp-proof-rescale}
\end{equation}

Euler's integral representation states that, for
\(\Re(c)>\Re(b)>0\),
\begin{equation}
  {}_2F_1(a_0,b;c;z)
  =
  \frac{\Gamma(c)}
       {\Gamma(b)\Gamma(c-b)}
  \int_0^1
  u^{b-1}(1-u)^{c-b-1}(1-zu)^{-a_0}\,\dd u.
  \label{eq:euler-hypergeometric}
\end{equation}
Apply it with \(a_0=1\), \(b=a\) and \(c=a+1\).  Since
\(\Gamma(a+1)=a\Gamma(a)\) and \(\Gamma(1)=1\),
\begin{equation}
  \int_0^1\frac{u^{a-1}}{1-zu}\,\dd u
  =\frac1a\,{}_2F_1(1,a;a+1;z).
  \label{eq:euler-specialised}
\end{equation}
Substituting this into \cref{eq:hyp-proof-rescale}, and using
\(ka=h\), gives
\begin{equation}
  \frac{v^a}{ph}
  {}_2F_1\!\left(1,\frac hk;\frac hk+1;-\frac qp v\right).
\end{equation}
Finally \(v^a=\e^{h\tau}\), which proves
\cref{eq:hypergeometric-integral} on the Euler domain.  Both sides are analytic
antiderivatives away from the singular set, so continuation along any admissible
interval preserves the identity.
\end{proof}

\subsection*{Application to the characteristic}

For \(\beta\neq\mu\), write the interior characteristic as
\begin{equation}
  x(\tau)
  =1+(1-\rho)
    \frac{Z\e^{b\tau}}{1-Z\e^{b\tau}},
  \qquad
  \rho=\frac{\mu}{\beta},
  \qquad b=\beta-\mu.
  \label{eq:appendix-x-characteristic}
\end{equation}
The first term integrates elementarily in
\(\alpha\int\e^{\alpha\tau}x(\tau)\,\dd\tau\).  Away from
\(\alpha+b=0\), the second is an instance of
\cref{prop:hypergeometric-integral} after taking
\(h=\alpha+b\), \(k=b\), \(p=1\) and \(q=-Z\), together with a factor
\(\alpha(1-\rho)Z\).  The Euler-domain calculation applies when its stated
hypotheses hold; elsewhere the identity is continued along the regular
characteristic.  Algebraic simplification by
\cref{eq:pochhammer-ratio} returns
\begin{equation}
  \alpha\int\e^{\alpha\tau}x(\tau)\,\dd\tau
  =\e^{\alpha\tau}\Psi(Z\e^{b\tau})+C.
\end{equation}
This agrees with \cref{eq:Psi-antiderivative}.  The series proof in the main text
has the advantage that the cancellation
\(\alpha+nb=b(n+\delta)\) remains visible; the integral identity has the advantage
that it applies to other exponential--rational characteristics without rebuilding
the calculation term by term.  At \(\alpha+b=0\), the same conclusion follows by
the finite resonant limit described after \cref{eq:abd-resonance}.
